\documentclass[11pt]{combiArticle}
\usepackage{tikz,amsmath,amsfonts,latexsym}
\usepackage{amssymb}
\usetikzlibrary{shapes,arrows,calc}
\usepackage{calc}
\usepackage{parskip}
\usepackage{array}
\setlength{\emergencystretch}{2em}

\definecolor{myblue}{RGB}{0, 84, 159}

\lang{de}

% Abkuerzungen fuer die Notation
\newcommand{\Dm}{D^{-}}
\newcommand{\Db}{\bar{D}}
\newcommand{\Dp}{D^{+}}
\newcommand{\Wr}{\mathcal{W}}
\newcommand{\ind}[1]{\mathbf{1}\!\left[#1\right]}
\newcommand{\code}[1]{\texttt{#1}}

\begin{document}
\course{Schichtplan Manager}
\assistent{Linus Palm}
\contact{palm@combi.rwth-aachen.de}
\title{Modell}
\date{09. November 2026}
\maketitle

Dieses Dokument beschreibt das gemischt-ganzzahlige lineare Programm (MILP), mit dem der Gurobi-Solver
(\code{api/app/solver/milp/service.py}, Klasse \code{GurobiSolverService}) den Dienstplan eines Monats
erstellt. Eine Fassung ohne Formeln ist die \emph{Modell -- Kurzfassung}.

%==========================================================================================================
\section{Überblick}

\textbf{Was entschieden wird.} Für jede Person, jeden Tag des Monats und jede Schicht: arbeitet sie diese Schicht
und auf welcher Station? Das Modell ist allgemein für mehrere Stationen formuliert; die Demo plant genau eine,
die fiktive Musterstation (\code{api/app/demo.py}).

\textbf{Feste Regeln und Ziele.} Jeder Plan hält die festen Regeln aus Abschnitt~\ref{sec:regeln} ein. Alles
andere ist ein Ziel. Die elf Ziele werden lexikografisch optimiert (Abschnitt~\ref{sec:lex}), also nacheinander
in dieser Rangfolge, wobei ein späteres Ziel ein früheres nicht verschlechtern darf:

\begin{center}
\small
\begin{tabular}{r l l l}
\hline
Rang & Ziel & Code & Abschnitt \\
\hline
1 & Monatssaldo innerhalb $\pm 4$\,h & \code{balance\_violation} & \ref{sec:z:saldo} \\
2 & Wochenende über den Monatswechsel fortführen & \code{weekend\_start} & \ref{sec:z:wochenende} \\
3 & möglichst wenige Lücken & \code{gaps} & \ref{sec:z:luecken} \\
4 & Lücken gleichmäßig auf die Tage verteilen & \code{gap\_spread} & \ref{sec:z:lueckenverteilung} \\
5 & Tage gleichmäßig besetzen & \code{staffing\_spread} & \ref{sec:z:tage} \\
6 & Samstage je Qualifikation gleichmäßig verteilen & \code{saturdays} & \ref{sec:z:samstage} \\
7 & Wünsche erfüllen & \code{wishes} & \ref{sec:z:wuensche} \\
8 & Rotation: jede Schichtart in jedem 28-Tage-Fenster & \code{rotation} & \ref{sec:z:rotation} \\
9 & Schichten eines Tages gleichmäßig besetzen & \code{shift\_balance} & \ref{sec:z:schichten} \\
10 & möglichst wenige Wechsel Spät $\to$ Früh & \code{Min\_L\_to\_E} & \ref{sec:z:spaetfrueh} \\
11 & möglichst oft dieselbe Schicht wie am Vortag & \code{cons} & \ref{sec:z:folgen} \\
\hline
\end{tabular}
\end{center}

\textbf{Monatswechsel.} Zwei Dinge verbinden aufeinanderfolgende Monate: Mit geladenem Vormonat zählen seine
Dienste als feste Werte bei allen Regeln mit. Endet ein Monat an einem Samstag, wird zudem der folgende Sonntag
mitgeplant (Abschnitt~\ref{sec:regeln}), damit jeder Samstagsdienst im nächsten Monat weitergehen kann.

\textbf{Lesehilfe.} Die Tage werden fortlaufend nummeriert: $1$ ist der erste Tag des Monats, $n$ der letzte,
$0$ der letzte Tag davor. Dienste vor dem Monat sind feste Werte, keine Variablen. Jeder Abschnitt gibt zuerst
die Formeln und darunter die Erklärung mit Verweis auf die Gleichungsnummern; in Klammern steht der Name im Code.

%==========================================================================================================
\section{Mengen und Indizes}

\begin{align}
  & i \in I && \text{Mitarbeitende} \label{set:I}\\
  & u \in U,\quad U^{St} \subseteq U && \text{Planungseinheiten, davon Stationen} \label{set:U}\\
  & s \in S && \text{Schichten} \label{set:S}\\
  & \theta \in \Theta = \{\text{Früh}, \text{Zwischen}, \text{Spät}, \text{Nacht}\},\quad S^{\theta} \subseteq S
    && \text{Schichtarten und ihre Schichten} \label{set:Theta}\\
  & l \in L = \{\text{Fachkraft}, \text{Hilfskraft}, \text{Azubi}, \text{MFA}\} && \text{Qualifikationen} \label{set:L}\\
  & d \in D = \{1,\dots,n\},\quad \Dp = D \cup \{n+1\} && \text{Tage des Monats, dazu der Folgetag} \label{set:D}\\
  & \hat D = \{1,\dots,\hat n\} && \text{Planungstage: der Monat, ggf.\ mit Sonntag} \label{set:Dhat}\\
  & d \in \Dm = \{1-m,\dots,0\},\quad \Db = \Dm \cup D && \text{Vortage mit bekannten Diensten, Planungshorizont} \label{set:Dm}\\
  & D^{Sa} \subseteq \{0,\dots,\hat n-1\} && \text{Samstage, deren Sonntag geplant wird} \label{set:DSa}\\
  & W = (w_1,\dots,w_{|W|}) \subseteq \{-m,\dots,n\} && \text{Samstage der Wochenenden mit einem Tag in } \Db \label{set:W}\\
  & D^{end} = \{1-\max(\Delta^N, L^N+1),\dots,\hat n-1\} && \text{Tage, an denen ein Nachtblock enden kann} \label{set:Dend}\\
  & K \subseteq S \times D \times U^{St} \times L && \text{Bedarfszeilen mit } r_{sdul} > 0 \label{set:K}\\
  & A \subseteq I \times S \times \hat D \times U^{St} && \text{zulässige Dienste, siehe \eqref{eq:A}} \label{set:A}
\end{align}

\textbf{Erklärung.}
\eqref{set:I}, \eqref{set:U}: $I$ (\code{I}) sind alle Mitarbeitenden der gewählten Stationen und ihrer
Springerpools, $U$ (\code{U}) alle Planungseinheiten. Dienste gibt es nur auf Stationen $U^{St}$
(\code{dataset.station\_ids}); Springerpools sind nur Herkunft.
\eqref{set:S}, \eqref{set:Theta}: Jede Schicht hat eine Art. Kurz geschrieben sind $S^{F}$, $S^{S}$, $S^{N}$ die
Früh-, Spät- und Nachtschichten (\code{S\_early}, \code{S\_late}, \code{S\_night}). Jede Regel gilt so für alle
Schichten einer Art, unabhängig vom Kürzel. In der Demo hat jede Art genau eine Schicht: F, Z, S und N.
\eqref{set:D}, \eqref{set:Dm}: $D$ (\code{D}) ist der Monat. $\Dm$ (\code{D\_prev}) sind die Vortage, für die
die tatsächlichen Dienste bekannt sind: Ist der Dienstplan des Vormonats geladen, ist $\Dm$ der ganze Vormonat;
sonst sind es die $m = 5$ Tage davor, an denen niemand gearbeitet hat. Der Folgetag $n+1$ in $\Dp$ zählt für
Abwesenheiten, in die eine Nacht vom letzten Monatstag hineinreicht.
\eqref{set:Dhat}: $\hat D$ (\code{D\_plan}) sind die Tage mit Dienstvariablen. Endet der Monat an einem Samstag,
gehört der Sonntag $n+1$ (\code{D\_next}) dazu und es ist $\hat n = n+1$, sonst $\hat n = n$; siehe
Abschnitt~\ref{sec:regeln}.
Allgemein gilt: Alle Tagesbereiche, die in die Vortage reichen (nicht-positive $d$), sind mit $\Dm$ geschnitten,
umfassen also nur bekannte Vortage.
\eqref{set:DSa}, \eqref{set:W}: $D^{Sa}$ (\code{D\_s}) enthält die Samstage, deren Sonntag im Monat liegt oder
mitgeplant wird, dazu den letzten Vortag, wenn er ein Samstag ist;
$W$ (\code{W}) jedes Wochenende, dessen Samstag oder Sonntag im Monat oder in den Vortagen liegt, durch seinen
Samstag und aufsteigend. Damit gehört das letzte Wochenende vor dem Monat dazu, gleich an welchem Wochentag der
Monat beginnt.
\eqref{set:Dend}: $D^{end}$ (\code{D\_end}) sind die Tage, für die ein Blockende in den Nachtregeln gebraucht
wird.
\eqref{set:K}: $K$ (Schlüssel von \code{R\_min}) enthält nur Zeilen mit tatsächlichem Bedarf.

%==========================================================================================================
\section{Parameter}

\subsection{Eingabedaten des Monats}

\begin{align}
  & r_{sdul} \in \mathbb{N} && \text{Mindestbesetzung von Schicht } s \text{ am Tag } d \text{ auf Station } u \text{ mit Qualifikation } l \label{par:r}\\
  & h_s \in \mathbb{N} && \text{bezahlte Minuten der Schicht } s \label{par:h}\\
  & t_i,\ c_i \in \mathbb{N} && \text{Sollminuten und gutgeschriebene Minuten im Monat} \label{par:t}\\
  & \lambda_{iud} \in L && \text{Qualifikation der an Tag } d \text{ gültigen Zugehörigkeit von } i \text{ zu } u \label{par:lambda}\\
  & y^{0}_{isd} \in \{0,1\},\ d \in \Dm && \text{fester Dienst } s \text{ von } i \text{ an Vortag } d \label{par:y0}\\
  & T_i \subseteq \Dp && \text{Fortbildungstage von } i \label{par:T}\\
  & B \subseteq I \times \Dp,\ V \subseteq B && \text{blockierte Tage, davon Urlaubstage} \label{par:B}\\
  & \operatorname{home}(i,d) \in U && \text{Stammeinheit von } i \text{ am Tag } d \label{par:home}
\end{align}

Abgeleitet aus dem Bedarf:
\begin{align}
  & R_{usd} = \textstyle\sum_{l} r_{sdul},\quad R_{ud} = \sum_{s} R_{usd},\quad R_{u} = \sum_{d \in D} R_{ud} \label{par:R}\\
  & U^{D} = \{u \in U^{St} : R_u > 0\} \label{par:UD}\\
  & \mathcal{B} = \{(u,s,l) : \exists\, d \in D \text{ mit } r_{sdul} > 0\} \label{par:demanded}\\
  & I_{udl} = \{i \in I : \lambda_{iud} = l\} \label{par:Iudl}
\end{align}

\textbf{Erklärung.}
\eqref{par:r} (\code{R\_min}) stammt aus der Seite „Mindestbesetzung“.
\eqref{par:h} (\code{h}) sind die Netto-Minuten der Schicht, die auf das Monatskonto gebucht werden. In der
Demo: F 05:55--13:25 (420\,min), Z 08:30--14:15 (345\,min), S 13:15--21:00 (435\,min),
N 20:10--06:10 (555\,min).
\eqref{par:t} (\code{t}, \code{credits}): In der Demo ist das Tagessoll einer Vollzeitkraft 462\,min
(38,5\,h je Woche) an jedem Werktag von Montag bis Freitag, der kein Feiertag ist; Teilzeitkräfte haben es
anteilig. Für Urlaub und Fortbildung an diesen Tagen wird das Tagessoll gutgeschrieben.
\eqref{par:lambda} (\code{level\_at}): Ein Dienst zählt mit der Qualifikation der Stationszugehörigkeit,
nicht mit der der Person.
\eqref{par:y0} (\code{past\_shifts}) sind die Dienste des geladenen Vormonats; ohne Vormonat sind alle 0.
\eqref{par:T} (\code{training}, \code{training\_days}), \eqref{par:B} (\code{blocked}, \code{vacation}): $B$
enthält alle Tage mit Urlaub, freiem Tag, Fortbildung oder „nicht verfügbar“, $V$ nur die Urlaubstage. Alle drei
umfassen auch den Folgetag $n+1$ (\code{context.availability}).
\eqref{par:home} (\code{home\_of}) ist die Einheit der an $d$ gültigen Stammzugehörigkeit.
\eqref{par:R}--\eqref{par:Iudl}: $R_{usd}$, $R_{ud}$, $R_u$ (\code{required\_usd}, \code{required\_ud},
\code{required\_u}) sind der Bedarf je Schicht, Tag und Monat über alle Qualifikationen. $\mathcal{B}$
(\code{demanded}) enthält die Kombinationen aus Station, Schicht und Qualifikation, die im Monat überhaupt
Bedarf haben. $I_{udl}$ (\code{staff\_for}) sind die Personen, die an $d$ auf $u$ als $l$ zählen.

\subsection{Im Projekt gesetzte Parameter}

\begin{center}
\small
\begin{tabular}{l l l >{\raggedright\arraybackslash}p{5.2cm}}
\hline
Symbol & Code & Wert & Bedeutung \\
\hline
$\tau$ & \code{POLICY.balance\_tolerance\_minutes} & 240 & erlaubte Saldoabweichung in Minuten (4\,h) \\
$K^{S}$ & \code{POLICY.max\_consecutive\_work\_days} & 10 & höchstens so viele Arbeitstage in Folge \\
$N^{max}$ & \code{POLICY.max\_consecutive\_nights} & 3 & höchstens so viele Nächte in Folge \\
$L^{N}$ & \code{L\_n} & $\lfloor 2880/1440 \rfloor = 2$ & ganz freie Tage nach einem Nachtblock \\
$\Delta^{N}$ & \code{Dist\_n} & 7 & Tage zwischen Ende eines Nachtblocks und Beginn des nächsten \\
$\Delta^{R}$ & \code{Rot\_w} & 28 & Tage eines Rotationsfensters \\
$\pi$ & \code{transfer\_penalty} & 10 & Strafe für einen Dienst auf fremder Station \\
$\omega$ & \code{wish\_reward} & 30 & Gewicht eines Wunsches \\
\hline
\end{tabular}
\end{center}

%==========================================================================================================
\section{Zulässige Dienste}

\begin{align}
  N^{V} &= \{(i,s,d) : \text{für } (i,d) \text{ ist eine Einschränkung eingetragen, die } s \text{ nicht erlaubt}\} \label{eq:NV}\\
  N^{N} &= \{(i,s,d) : s \in S^{N},\ d \in D,\ (i,d+1) \in B\} \label{eq:NN}\\
  N^{S} &= \{(i,s,d) : s \in S^{S},\ d \in D,\ \bigl((i,d+1) \in V,\ (i,d) \notin V\bigr) \text{ oder }
          \bigl(d+1 \in T_i,\ d \notin T_i\bigr)\} \label{eq:NS}\\
  A &= \bigl\{(i,s,d,u) : u \in U^{St},\ \lambda_{iud} \text{ definiert},\
       (i,s,d) \notin N^{V} \cup N^{N} \cup N^{S}, \label{eq:A}\\
    &\hphantom{{}= \bigl\{(i,s,d,u) :{}} s \in S^{F} \cup S^{S} \text{ oder } (u,s,\lambda_{iud}) \in \mathcal{B}\bigr\} \notag
\end{align}

\textbf{Erklärung.}
Statt diese Regeln als Nebenbedingungen zu formulieren, werden verbotene Dienste gar nicht erst als Variable
angelegt (\code{not\_allowed}, Liste beim Anlegen von \code{assignment}).
\eqref{eq:NV}: Urlaub, freier Tag, Fortbildung und „nicht verfügbar“ verbieten alle Schichten, „nur
bestimmte Schichten“ verbietet die übrigen. Für einen mitgeplanten Sonntag gelten die Einträge, die der Folgemonat
für ihn schon hat.
\eqref{eq:NN}: Eine Nacht reicht in den Folgetag; vor einem blockierten Tag ist sie verboten, am letzten
Monatstag also auch vor einer Abwesenheit am Ersten des Folgemonats.
\eqref{eq:NS}: Am Tag vor Urlaubsbeginn und vor einer Fortbildung gibt es keinen Spätdienst, auch wenn sie am
Ersten des Folgemonats beginnen. Vor dem Urlaub ist das betriebliche Gepflogenheit; die Fortbildung ist
Arbeitszeit, und nach einem Spätdienst bis 21:00 blieben 11\,h Ruhezeit nur bis zu einem Beginn um 08:00, dessen
Uhrzeit das Modell nicht kennt.
\eqref{eq:A} (\code{always\_open}, \code{demanded}): Ein Dienst ist nur auf einer Station möglich, zu der an
diesem Tag eine gültige Zugehörigkeit besteht. Früh- und Spätdienste kann dort jede Qualifikation übernehmen, auch
ohne eigenen Bedarf; sie zählen dann als zusätzliches Personal. Nacht- und Zwischendienste dagegen nur, wenn die
Station sie von dieser Qualifikation im Monat überhaupt fordert (z.\,B. kein Zwischendienst, wenn nie einer
gefordert ist, keine Nächte für Azubis, wenn nie Azubi-Nachtbedarf besteht).

%==========================================================================================================
\section{Variablen und Hilfsausdrücke}

\subsection{Entscheidungsvariablen}

\begin{align}
  & x_{isdu} \in \{0,1\} && (i,s,d,u) \in A && \text{Dienst } (\code{assignment}) \label{var:x}\\
  & \psi_{iw\theta} \in \{0,1\} && (i,w,\theta) \in \Psi && \text{Schichtart im Rotationsfenster } (\code{rotation}) \label{var:psi}\\
  & g_{sdul} \ge 0 && (s,d,u,l) \in K && \text{Lücke } (\code{gap}) \label{var:g}\\
  & \eta_i \in \mathbb{R} && i \in I && \text{Monatssaldo in Minuten } (\code{hour}) \label{var:eta}\\
  & o_i \ge 0 && i \in I && \text{Saldo außerhalb des Bands } (\code{over}) \label{var:o}\\
  & q_{ud} \ge 0 && u \in U^{St},\ d \in D,\ R_{ud} > 0 && \text{Quadrat der Lücken des Tages } (\code{gap\_sq}) \label{var:q}\\
  & \rho_{ud} \ge 0 && u \in U^{D},\ d \in D && \text{Abweichung der Tagesbesetzung } (\code{day\_dev}) \label{var:rho}\\
  & \chi_i \ge 0 && i \in I && \text{Quadrat der gearbeiteten Samstage } (\code{sat\_sq}) \label{var:chi}\\
  & \delta_{usd} \ge 0 && u \in U^{St},\ d \in D,\ R_{ud} > 0,\ s \in S && \text{Abweichung der Schichtbesetzung } (\code{dev}) \label{var:delta}\\
  & b_{id} \in [0,1] && i \in I,\ d \in \{0\} \cup D && \text{Beginn eines Nachtblocks } (\code{n\_beg}) \label{var:b}\\
  & e_{id} \in [0,1] && i \in I,\ d \in D^{end} && \text{Ende eines Nachtblocks } (\code{n\_end}) \label{var:e}\\
  & z_{id} \in [0,1] && i \in I,\ d \in D && \text{Spät- vor Frühdienst } (\code{L\_to\_E}) \label{var:z}\\
  & v_{isd} \in [0,1] && i \in I,\ s \in S,\ d \in D && \text{gleiche Schicht wie am Vortag } (\code{cons\_s}) \label{var:v}
\end{align}

\textbf{Erklärung.}
Die eigentliche Entscheidung ist $x$ \eqref{var:x}. $\psi$ \eqref{var:psi} ist ebenfalls binär; die Menge $\Psi$
steht bei der Rotation in Abschnitt~\ref{sec:z:rotation}. Alle anderen Variablen sind stetige Hilfsgrößen; ihre
Werte ergeben sich aus $x$ und der Optimierung. $g$ \eqref{var:g} nimmt im Optimum ganzzahlige Werte an, $b$, $e$,
$z$ und $v$ \eqref{var:b}--\eqref{var:v} nur 0 oder 1, weil sie ausschließlich durch ganzzahlige Ausdrücke in $x$
nach oben bzw. unten begrenzt werden und die Optimierung sie an diese Grenze drückt.

\subsection{Hilfsausdrücke}

\begin{align}
  X_{isd} &= \begin{cases}
      \sum_{u : (i,s,d,u) \in A} x_{isdu} & d \in \hat D\\
      y^{0}_{isd} & d \in \Dm\\
      0 & \text{sonst}
    \end{cases} && \text{(\code{x}, \code{shift\_on})} \label{eq:X}\\
  Y_{id} &= \textstyle\sum_{s \in S} X_{isd} && \text{(\code{work})} \label{eq:Y}\\
  Y^{F}_{id} &= \textstyle\sum_{s \in S^{F}} X_{isd},\quad
  Y^{S}_{id} = \sum_{s \in S^{S}} X_{isd},\quad
  Y^{N}_{id} = \sum_{s \in S^{N}} X_{isd} && \text{(\code{early}, \code{late}, \code{night})} \notag\\
  P_{usd} &= \textstyle\sum_{i : (i,s,d,u) \in A} x_{isdu} && \text{(\code{staffed})} \label{eq:P}\\
  P_{ud} &= \textstyle\sum_{s} P_{usd},\quad
  P_{u} = \sum_{d \in D} P_{ud} && \text{(\code{staff\_of\_day}, \code{staff\_of\_month})} \notag\\
  G_{ud} &= \textstyle\sum_{(s,d,u,l) \in K} g_{sdul} && \text{(\code{gaps\_of\_day})} \label{eq:G}
\end{align}

\textbf{Erklärung.}
\eqref{eq:X}: $X_{isd}$ ist 1, wenn $i$ an $d$ Schicht $s$ arbeitet, egal auf welcher Station. Auf den
Vortagen ist es der feste Dienst aus dem Vormonat, außerhalb des bekannten Zeitraums 0. Dadurch wirken alle
Regeln über den Monatswechsel, ohne dass die Vergangenheit verändert werden kann.
\eqref{eq:Y}: $Y_{id}$ ist 1, wenn $i$ an $d$ irgendeinen Dienst hat, $Y^{F}$, $Y^{S}$, $Y^{N}$ entsprechend
für Früh-, Spät- und Nachtdienst.
\eqref{eq:P}: Besetzung einer Station je Schicht, je Tag und im Monat.
\eqref{eq:G}: Summe der Lücken einer Station an einem Tag.

%==========================================================================================================
\section{Feste Regeln}\label{sec:regeln}

Diese Nebenbedingungen hält jeder Plan ein. Zusammen mit den verbotenen Diensten aus \eqref{eq:A} sind das alle
festen Regeln. Am Monatsanfang erzwingen sie nie etwas, das die Vergangenheit unmöglich machen könnte: So bleibt
das Modell auch mit geladenem Vormonat immer lösbar.

\textbf{Der mitgeplante Sonntag.} Die Regeln gelten auf allen Planungstagen $\hat D$. Endet der Monat an einem
Samstag, ist das auch der Sonntag $n+1$ (\code{D\_next}). Er ist durch \eqref{con:weekend} eine Kopie des Samstags:
dieselbe Schicht oder beide frei. So plant der Monat nur Samstagsdienste, die am Sonntag weitergehen können, ohne
eine vierte Nacht, gegen die Erholung oder eine Abwesenheit am Sonntag. Der Sonntag zählt aber für kein Konto
\eqref{con:balance}, keinen Bedarf und kein Ziel, und seine Dienste gehören nicht zum Plan. Der Folgemonat plant ihn
mit diesem Plan als Vormonat neu und führt das Wochenende über \eqref{obj:weekend} fort.

\subsection{Ein Dienst pro Tag und Arbeitstage am Stück}

\begin{align}
  & Y_{id} \le 1 && \forall i \in I,\ d \in \hat D \label{con:one}\\
  & \textstyle\sum_{d' = d}^{d + K^{S}} Y_{id'} + |T_i \cap \{d,\dots,d+K^{S}\}| \le \max\bigl(K^{S},\ \kappa_{id}\bigr)
    && \forall i \in I,\ d \in \{1-\min(m,K^{S}),\dots,\hat n-K^{S}\} \label{con:ks}
\end{align}

\textbf{Erklärung.}
\eqref{con:one} (\code{one\_shift}): Höchstens ein Dienst pro Person und Tag, über alle Stationen.
\eqref{con:ks}: In jedem Fenster von $K^{S}+1 = 11$ Tagen höchstens $K^{S} = 10$ Arbeitstage; Fortbildungstage
zählen als Arbeitstage. Die Fenster beginnen schon in den Vortagen, damit die Regel über den Monatswechsel gilt.
$\kappa_{id}$ ist der feste Teil des Fensters (Arbeitstage der Vortage und Fortbildungstage). Erreicht er die
Grenze schon allein, ist im Fenster kein weiterer Dienst möglich, statt dass das Modell unlösbar wird.

\subsection{Nachtdienste}

\begin{align}
  & \textstyle\sum_{d' = d}^{d + N^{max}} Y^{N}_{id'} \le \max\bigl(N^{max},\ \nu_{id}\bigr)
    && \forall i \in I,\ d \in \{1-\min(m,N^{max}),\dots,0\} \cup \{1,\dots,\hat n-N^{max}\} \label{con:nights}\\
  & b_{id} \ge Y^{N}_{id} - Y^{N}_{i,d-1} && \forall i \in I,\ d \in \{0\} \cup \hat D \label{con:beg}\\
  & e_{id} \ge Y^{N}_{id} - Y^{N}_{i,d+1} && \forall i \in I,\ d \in D^{end} \label{con:end}\\
  & b_{id} + e_{id} \le 1 && \forall i \in I,\ d \in \{1,\dots,\hat n-1\} \label{con:single}\\
  & e_{id} + Y_{i,d+j} \le 1 && \forall i \in I,\ d \in \{-L^{N},\dots,\hat n-1\},\ j = 1,\dots,L^{N},\ d+j \in \hat D \label{con:recovery}\\
  & e_{id} \le 0 && \forall i \in I,\ d \in \{1,\dots,\hat n-1\} \text{ mit } T_i \cap \{d+1,\dots,d+L^{N}\} \neq \emptyset
    \label{con:recoverytraining}\\
  & e_{id} + Y^{F}_{i,d+L^{N}+1} \le 1 && \forall i \in I,\ d \in \{-L^{N},\dots,\hat n-1\},\ d+L^{N}+1 \in \hat D \label{con:recoveryearly}\\
  & e_{id} + \textstyle\sum_{j = 1,\ d+j \in \hat D}^{\Delta^{N}} b_{i,d+j} \le 1 && \forall i \in I,\ d \in \{1-\Delta^{N},\dots,\hat n-1\} \cap D^{end} \label{con:distance}
\end{align}

\textbf{Erklärung.}
\eqref{con:nights}: In jedem Fenster von $N^{max}+1 = 4$ Tagen höchstens $N^{max} = 3$ Nächte; $\nu_{id}$ sind
die Nächte der Vortage im Fenster (wie bei \eqref{con:ks}).
\eqref{con:beg}, \eqref{con:end} (\code{n\_beg}, \code{n\_end}): $b_{id}$ ist mindestens 1, wenn an $d$ ein
Nachtblock beginnt (Nacht an $d$, keine an $d-1$), $e_{id}$, wenn einer endet. Untergrenzen genügen, weil $b$
und $e$ nur in $\le$-Bedingungen mit positivem Vorzeichen und in keinem Ziel vorkommen.
\eqref{con:single}: Beginn und Ende am selben Tag hieße Ein-Nacht-Block; der ist im Monat verboten. Eine
einzelne Nacht am letzten Vortag wird nicht zwangsweise fortgesetzt, weil eine Abwesenheit am Ersten oder die
Nachtgrenze die zweite Nacht verbieten kann.
\eqref{con:recovery}, \eqref{con:recoverytraining}, \eqref{con:recoveryearly}: Nach dem letzten Nachtdienst eines
Blocks sind $L^{N} = 2$ Tage ganz frei, und am dritten Tag gibt es keinen Frühdienst. Das bildet die 48 Stunden
Erholung ab: Die Nacht endet um 06:10, der Frühdienst am dritten Tag beginnt um 05:55. \emph{Ganz frei} heißt auch
ohne Fortbildung \eqref{con:recoverytraining} (\code{recovery\_training}): Ein Block im Monat endet nicht innerhalb
von $L^{N}$ Tagen vor einer Fortbildung, auch nicht vor einer am Folgetag $n+1$. Mit \eqref{eq:NN} gibt es so in
den $L^{N}$ Tagen vor einer Fortbildung gar keine Nacht. Nur Tage im Monat werden eingeschränkt; ein Block, der
schon vor dem Monat geendet hat, bleibt, wie er ist.
\eqref{con:distance}: Nach dem Ende eines Nachtblocks darf in den nächsten $\Delta^{N} = 7$ Tagen kein neuer
Block beginnen; berücksichtigt werden nur Blockanfänge im Monat.

\subsection{Wochenenden}

\begin{align}
  & X_{isd} = X_{is,d+1} && \forall i \in I,\ s \in S,\ d \in D^{Sa} \cap D \label{con:weekend}\\
  & X_{is1} \le X_{is0} && \forall i \in I,\ s \in S, \text{ falls } 0 \in D^{Sa},\ Y_{i0} = 1 \label{con:weekendstart}\\
  & \Omega_{iw} = \begin{cases}
      \max\bigl(Y_{iw},\, Y_{i,w+1}\bigr) & w \le 0\\
      Y_{iw} & w \ge 1
    \end{cases} && \text{(\code{weekend})} \label{eq:Omega}\\
  & \Omega_{i,w_k} + \Omega_{i,w_{k+1}} \le 1 && \forall i \in I,\ k = 1,\dots,|W|-1 \text{ mit } w_{k+1} \ge 0
    \label{con:twoweekends}
\end{align}

\textbf{Erklärung.}
\eqref{con:weekend}: Samstag und Sonntag haben dieselbe Schicht, oder beide sind frei.
\eqref{con:weekendstart}: Über den Monatswechsel (der Monat beginnt an einem Sonntag) bindet der bekannte Samstag
den Sonntag als feste Regel nur schwach: Nach einem gearbeiteten Samstag hat der Sonntag dieselbe Schicht oder ist
frei; erzwungen wird nichts, weil eine Abwesenheit am Ersten oder die Nachtgrenze den Dienst verbieten kann. War
der Samstag frei, ist der Sonntag offen; ohne Vormonat gelten die Vortage ja nur als frei, sonst müsste jeder am
Ersten frei haben. Mit geladenem Vormonat verlangt das Ziel \eqref{obj:weekend} die volle Regel.
\eqref{eq:Omega}: $\Omega_{iw}$ gibt an, ob $i$ das Wochenende mit dem Samstag $w$ arbeitet. Gearbeitet ist es
mit einem Dienst am Samstag oder Sonntag; eine Nacht vom Freitag allein lässt es frei. Im Monat haben Samstag
und Sonntag dieselbe Schicht \eqref{con:weekend}, also entscheidet der Samstag ($w \ge 1$). Für $w \le 0$ liegt
der Samstag vor dem Monat und $Y_{iw}$ fest (außerhalb der bekannten Tage 0); das Maximum ist dann linear: 1 nach
einem gearbeiteten Samstag, sonst $Y_{i,w+1}$.
\eqref{con:twoweekends}: Keine zwei Wochenenden in Folge, auch nicht über den Monatswechsel. Geprüft werden nur
Paare, deren späteres Wochenende den Monat berührt ($w_{k+1} \ge 0$, sein Sonntag liegt also im Monat); die
Vergangenheit wird nie gegen sich selbst geprüft, auch nicht ein Paar, das nach einem gearbeiteten Samstag vor
dem Monat schon ganz feststeht. Daraus folgt zugleich, dass höchstens jedes zweite Wochenende gearbeitet wird.

%==========================================================================================================
\section{Ziele in ihrer Rangfolge}\label{sec:ziele}

Jedes Ziel steht hier mit den Nebenbedingungen, die seine Hilfsvariablen festlegen. Die Reihenfolge ist die
Rangfolge, in der Gurobi sie optimiert (Abschnitt~\ref{sec:lex}).

\subsection{Rang 1: Monatssaldo innerhalb des Bands}\label{sec:z:saldo}

\begin{align}
  & \textstyle\sum_{(s,d,u) : (i,s,d,u) \in A,\ d \in D} h_s\, x_{isdu} + c_i - \eta_i = t_i && \forall i \in I \label{con:balance}\\
  & o_i \ge \eta_i - \tau,\qquad o_i \ge -\eta_i - \tau && \forall i \in I \label{con:over}\\
  & f_1 = \textstyle\sum_{i \in I} o_i \to \min && \text{(\code{balance\_violation})} \label{obj:balance}
\end{align}

\textbf{Erklärung.}
\eqref{con:balance} (\code{workhours}): $\eta_i$ ist der Saldo in Minuten: geplante Minuten im Monat plus
Gutschriften minus Soll, genau wie das Monatskonto im Projekt. Ein mitgeplanter Sonntag zählt nicht mit.
\eqref{con:over} (\code{over\_up}, \code{over\_down}): $o_i$ misst, um wie viele Minuten der Saldo außerhalb von
$\pm\tau = \pm 240$ liegt; im Optimum ist $o_i = \max(0, |\eta_i| - \tau)$.
\eqref{obj:balance}: Dieses Ziel steht ganz oben, noch vor den Lücken. Können alle Salden im Band liegen, ist
$\hat f_1 = 0$, und jede spätere Stufe muss $f_1 = 0$ halten \eqref{obj:lex}: Das Band wirkt dann wie eine feste
Regel, auch wenn dafür eine Lücke offen bleibt. Ist es für jemanden unerreichbar (z.\,B.\ den halben Monat nicht
verfügbar), entsteht trotzdem ein Plan mit möglichst kleiner Überschreitung, statt dass das Modell unlösbar wird.

\subsection{Rang 2: Wochenende über den Monatswechsel}\label{sec:z:wochenende}

\begin{align}
  & f_2 = \textstyle\sum_{i \in I : Y_{i0} = 1} \bigl(1 - X_{i\sigma_i 1}\bigr) + \sum_{i \in I : Y_{i0} = 0} Y_{i1} \to \min
    && \text{(\code{weekend\_start})} \label{obj:weekend}
\end{align}
nur mit geladenem Vormonat und wenn der Monat an einem Sonntag beginnt; sonst ist $f_2 = 0$. Dabei ist $\sigma_i$ die
Schicht von $i$ am Samstag $0$.

\textbf{Erklärung.}
\eqref{obj:weekend} (\code{weekend\_split}): Zählt die Personen, deren Wochenende über den Monatswechsel ungleich
ist: Wer Samstag gearbeitet hat und am Sonntag nicht dieselbe Schicht hat, und wer Samstag frei hatte und am Sonntag
arbeitet. Andere Schichten als die vom Samstag schließt \eqref{con:weekendstart} schon aus. Weil das Ziel direkt
nach dem Saldo kommt, gilt die volle Wochenendregel, wann immer das möglich ist. Ist es nicht möglich, etwa weil
nach dem Plan des Vormonats noch Urlaub am Ersten eingetragen wurde, bleibt das Wochenende geteilt, statt dass das
Modell unlösbar wird; die Prüfung weist es dann aus. Hat der Vormonat mit seinem mitgeplanten Sonntag gerechnet,
kann jeder Samstagsdienst weitergehen. Ohne Vormonat gilt dieses Ziel nicht, weil die Tage davor nur als frei
angenommen sind.

\subsection{Rang 3: Lücken}\label{sec:z:luecken}

\begin{align}
  & \textstyle\sum_{i \in I_{udl} : (i,s,d,u) \in A} x_{isdu} + g_{sdul} \ge r_{sdul} && \forall (s,d,u,l) \in K \label{con:demand}\\
  & f_3 = \textstyle\sum_{(s,d,u,l) \in K} g_{sdul} \to \min && \text{(\code{gaps})} \label{obj:gaps}
\end{align}

\textbf{Erklärung.}
\eqref{con:demand} (\code{demand}): Die Mindestbesetzung jeder Bedarfszeile wird durch Dienste mit passender
Qualifikation gedeckt; was fehlt, ist eine Lücke $g$. Die Mindestbesetzung ist damit keine feste Regel: Auch bei
Personalmangel gibt es einen Plan, und \eqref{obj:gaps} hält die Zahl der unbesetzten Pflichtstellen klein.

\subsection{Rang 4: Lücken gleichmäßig verteilen}\label{sec:z:lueckenverteilung}

\begin{align}
  & q_{ud} \ge (2k+1)\, G_{ud} - k(k+1) && \forall (u,d) : R_{ud} > 0,\ k = 0,\dots,R_{ud}-1 \label{con:gapsq}\\
  & f_4 = \textstyle\sum_{(u,d) : R_{ud} > 0} \dfrac{q_{ud}}{R_{ud}} \to \min && \text{(\code{gap\_spread})} \label{obj:gapspread}
\end{align}

\textbf{Erklärung.}
\eqref{con:gapsq} (\code{gap\_sq}): Lineare Darstellung von $q_{ud} \ge G_{ud}^2$. Die rechte Seite ist für
jedes $k$ die Sekante durch die Parabelpunkte $(k, k^2)$ und $(k+1, (k+1)^2)$. Zwischen $k$ und $k+1$ liegt sie
über der Parabel, außerhalb darunter. Das Maximum aller Sekanten ist daher die stückweise lineare Interpolation
von $G_{ud}^2$: überall mindestens $G_{ud}^2$ und in ganzzahligen Punkten genau $G_{ud}^2$. Da die Lücken eines
Tages ganzzahlig zwischen $0$ und $R_{ud}$ liegen und $q_{ud}$ minimiert wird, ist $q_{ud} = G_{ud}^2$ im
Optimum exakt.
\eqref{obj:gapspread}: Durch das Quadrat sind zwei Tage mit je einer Lücke ($1+1$) billiger als ein Tag mit zwei
($4$). Die Division durch $R_{ud}$ gewichtet relativ zum Bedarf: Eine Lücke wiegt an einem schwach besetzten Tag
schwerer.

\subsection{Rang 5: Tage gleichmäßig besetzen}\label{sec:z:tage}

\begin{align}
  & \rho_{ud} \ge P_{ud} - \tfrac{R_{ud}}{R_u}\, P_u,\qquad
    \rho_{ud} \ge -\bigl(P_{ud} - \tfrac{R_{ud}}{R_u}\, P_u\bigr) && \forall u \in U^{D},\ d \in D \label{con:daydev}\\
  & f_5 = \textstyle\sum_{u \in U^{D}} \sum_{d \in D} \rho_{ud} \to \min && \text{(\code{staffing\_spread})} \label{obj:staffingspread}
\end{align}

\textbf{Erklärung.}
\eqref{con:daydev} (\code{day\_dev\_up}, \code{day\_dev\_down}): $\rho_{ud} \ge |P_{ud} - \frac{R_{ud}}{R_u}
P_u|$, also die Abweichung der Tagesbesetzung von ihrem Anteil an der Monatsbesetzung. Ein Tag ohne Bedarf
hat Anteil 0, jeder Dienst dort zählt als Abweichung.
\eqref{obj:staffingspread}: Jeder Tag einer Station soll im gleichen Verhältnis Ist/Soll besetzt sein wie der
ganze Monat; Über- und Unterbesetzung verteilen sich so gleichmäßig auf die Tage.

\subsection{Rang 6: Samstage gleichmäßig verteilen}\label{sec:z:samstage}

\begin{align}
  & a_i = \textstyle\sum_{d \in D,\ d \text{ Samstag}} Y_{id} && \forall i \in I \label{eq:sat}\\
  & \chi_i \ge (2k+1)\, a_i - k(k+1) && \forall i \in I,\ k = 0,\dots,s^{Sa}-1 \label{con:satsq}\\
  & f_6 = \textstyle\sum_{i \in I} \chi_i \to \min && \text{(\code{saturdays})} \label{obj:saturdays}
\end{align}

\textbf{Erklärung.}
\eqref{eq:sat} (\code{saturdays\_of}): $a_i$ ist die Zahl der Samstage im Monat, an denen $i$ arbeitet; $s^{Sa}$ ist
die Zahl der Samstage im Monat. Wegen \eqref{con:weekend} ist das zugleich die Zahl der gearbeiteten Wochenenden.
\eqref{con:satsq} (\code{sat\_sq}): $\chi_i \ge a_i^2$, mit denselben Sekanten wie bei \eqref{con:gapsq}; im Optimum
ist $\chi_i = a_i^2$.
\eqref{obj:saturdays}: Bei gegebener Zahl gearbeiteter Samstage ist die Summe der Quadrate am kleinsten, wenn sich die
Zahlen höchstens um eins unterscheiden: $2+2+3$ ist besser als $1+3+3$. Jede Qualifikation deckt ihren eigenen
Bedarf, deshalb gleicht die Summe über alle Personen die Samstage innerhalb jeder Qualifikation aus. Wer an einem
Samstag nicht kann (Urlaub, Einschränkung), hat entsprechend weniger; die übrigen teilen den Rest gleichmäßig. Die
Zahlen sind absolut, also unabhängig von Teilzeit. Durch die Sekanten ist die Schranke der Relaxation in der Regel
schon der ganzzahlige Wert, sodass Gurobi die Stufe schnell exakt löst.

\subsection{Rang 7: Wünsche}\label{sec:z:wuensche}

\begin{align}
  p_{isdu} = {} & - \pi \cdot \ind{\operatorname{home}(i,d) \in U^{St},\ \operatorname{home}(i,d) \neq u}
                 \label{par:p}\\
  & - \omega \cdot \ind{\text{Wunsch „freier Tag“ für } (i,d)}
    - \omega \cdot \ind{s \in S^{N},\ \text{Wunsch „freier Tag“ für } (i,d+1)} \notag\\
  & - \omega \cdot \ind{\text{Wunsch „Schicht } s \text{ frei“ für } (i,d)} \notag\\
  & + \omega \cdot \ind{\text{Wunsch „Wunschtag“ für } (i,d)}
    + \omega \cdot \ind{\text{Wunsch „Wunschschicht } s\text{“ für } (i,d)} \notag\\
  f_7 = {} & \textstyle\sum_{(i,s,d,u) \in A} p_{isdu}\, x_{isdu} \to \max \qquad \text{(\code{wishes})} \label{obj:wishes}
\end{align}

\textbf{Erklärung.}
\eqref{par:p} (\code{pref}): Jeder zulässige Dienst bekommt einen Präferenzwert. Ein erfüllter Dienstwunsch bringt
$+\omega$, ein Dienst gegen einen Freiwunsch kostet $\omega$. Beim freien Tag zählt auch ein Nachtdienst am
Vorabend, weil er bis 06:10 in den Wunschtag hineinreicht; so bewertet es auch die Prüfung. Ein Dienst auf einer
fremden Station kostet $\pi$, aber nur, wenn die Person eine Station als Stammeinheit hat; Springer (Stammeinheit
ist ein Pool) dürfen überall gleich gut arbeiten. In der Demo mit nur einer Station kommt das nicht vor.
\eqref{obj:wishes}: Die Summe der Präferenzwerte aller Dienste wird maximiert.

\subsection{Rang 8: Rotation}\label{sec:z:rotation}

\begin{align}
  & \Wr = \{w \in \Db : 1 \le w + \Delta^{R} - 1 \le n\},\qquad [w] = \{w,\dots,w + \Delta^{R} - 1\} \label{set:Wr}\\
  & \Psi = \bigl\{(i,w,\theta) : w \in \Wr,\ \ y^{0}_{isd} = 0\ \forall s \in S^{\theta},\ d \in [w] \cap \Dm,
    \label{set:Psi}\\
  & \hphantom{\Psi = \bigl\{(i,w,\theta) :{}}
    \exists\,(i,s,d,u) \in A \text{ mit } s \in S^{\theta},\ d \in [w] \cap D\bigr\} \notag\\
  & \psi_{iw\theta} \le \textstyle\sum_{d \in [w] \cap D} \sum_{s \in S^{\theta}} X_{isd} && \forall (i,w,\theta) \in \Psi \label{con:rotation}\\
  & f_8 = \textstyle\sum_{(i,w,\theta) \in \Psi} \psi_{iw\theta} \to \max && \text{(\code{rotation})} \label{obj:rotation}
\end{align}

\textbf{Erklärung.}
Niemand soll wochenlang nur dieselbe Schicht haben: In jedem Fenster von $\Delta^{R} = 28$ Tagen soll jede Person
möglichst jede Schichtart haben, die sie arbeiten kann.
\eqref{set:Wr} (\code{window\_starts}): $\Wr$ sind die ersten Tage aller Fenster, die im Monat enden und an
einem bekannten Tag beginnen. Mit geladenem Vormonat reichen sie bis zu 27 Tage in ihn zurück, sodass die
Rotation über den Monatswechsel weiterläuft; ohne ihn beginnen sie frühestens $m = 5$ Tage vor dem Monat (ein
Monat mit 31 Tagen hat dann 9 Fenster).
\eqref{set:Psi} (Schlüssel von \code{worked\_type}): Gezählt wird ein Paar aus Fenster und Schichtart nur, wenn die
Person diese Art im Fensterteil des Monats überhaupt arbeiten kann, also eine Variable dafür existiert. Früh und
Spät kann jede Qualifikation \eqref{eq:A}; Azubis werden aber nicht in Nächte gedrängt, wenn nie
Azubi-Nachtbedarf besteht, und „nur bestimmte Schichten“ lässt die übrigen Arten aus. Hat die Person die Art im Fensterteil vor dem Monat schon gearbeitet, ist das Paar erfüllt und
braucht keine Variable.
\eqref{con:rotation}: $\psi_{iw\theta}$ kann nur 1 sein, wenn $i$ im Fenster ab $w$ mindestens eine Schicht der Art
$\theta$ arbeitet.
\eqref{obj:rotation}: Die Zahl der erfüllten Paare aus Person, Fenster und Schichtart wird maximiert. Ohne
Vormonat überdecken alle Fenster die Monatsmitte; dann genügt eine Schicht jeder Art dort für alle Fenster.

\subsection{Rang 9: Schichten eines Tages gleichmäßig besetzen}\label{sec:z:schichten}

\begin{align}
  & \delta_{usd} \ge P_{usd} - \tfrac{R_{usd}}{R_{ud}}\, P_{ud},\qquad
    \delta_{usd} \ge -\bigl(P_{usd} - \tfrac{R_{usd}}{R_{ud}}\, P_{ud}\bigr) && \forall (u,d) : R_{ud} > 0,\ s \in S \label{con:shiftdev}\\
  & f_9 = \textstyle\sum_{u,s,d} \delta_{usd} \to \min && \text{(\code{shift\_balance})} \label{obj:shiftbalance}
\end{align}

\textbf{Erklärung.}
\eqref{con:shiftdev} (\code{shift\_dev\_up}, \code{shift\_dev\_down}): $\delta_{usd} \ge |P_{usd} -
\frac{R_{usd}}{R_{ud}} P_{ud}|$, die Abweichung einer Schicht von ihrem Anteil an der Tagesbesetzung.
\eqref{obj:shiftbalance}: Zusätzliches Personal eines Tages verteilt sich im Verhältnis der Mindestbesetzung auf
die Schichten. Weil sich Anteile nur mit „Bruchteilen von Personen“ exakt treffen lassen, ist ein Teil dieser
Abweichung unvermeidbar; die Stufe hat deshalb ein Zeitlimit statt eines Optimalitätsbeweises.

\subsection{Rang 10: Wechsel vom Spät- in den Frühdienst}\label{sec:z:spaetfrueh}

\begin{align}
  & z_{id} \ge Y^{S}_{i,d-1} + Y^{F}_{id} - 1 && \forall i \in I,\ d \in D \label{con:ltoe}\\
  & f_{10} = \textstyle\sum_{i \in I} \sum_{d \in D} z_{id} \to \min && \text{(\code{Min\_L\_to\_E})} \label{obj:ltoe}
\end{align}

\textbf{Erklärung.}
\eqref{con:ltoe} (\code{L\_to\_E}): $z_{id}$ ist mindestens 1, wenn auf einen Spätdienst am Vortag ein Frühdienst
folgt, auch über den Monatswechsel.
\eqref{obj:ltoe}: Zwischen Spätdienst-Ende (21:00) und Frühdienst-Beginn (05:55) liegen nur 8:55\,h. Mit den
Schichtzeiten der Demo und den Nachtregeln ist das der einzige mögliche Fall unter 11\,h Ruhezeit. Er ist bewusst
erlaubt, wird aber so selten wie möglich geplant.

\subsection{Rang 11: Dieselbe Schicht wie am Vortag}\label{sec:z:folgen}

\begin{align}
  & v_{isd} \le X_{isd},\qquad v_{isd} \le X_{is,d-1} && \forall i \in I,\ s \in S,\ d \in D \label{con:cons}\\
  & f_{11} = \textstyle\sum_{i \in I} \sum_{s \in S} \sum_{d \in D} v_{isd} \to \max && \text{(\code{cons})} \label{obj:cons}
\end{align}

\textbf{Erklärung.}
\eqref{con:cons} (\code{cons\_s}): $v_{isd}$ darf nur 1 sein, wenn an $d$ und $d-1$ dieselbe Schicht gearbeitet
wird.
\eqref{obj:cons}: Möglichst oft dieselbe Schicht wie am Vortag, damit ruhige Folgen entstehen.

%==========================================================================================================
\section{Lexikografisches Lösen}\label{sec:lex}

Für $k = 1,\dots,11$ wird $f_k$ über alle zulässigen Lösungen optimiert, unter den zusätzlichen Bedingungen
\begin{equation}
  f_j \le \hat f_j + \alpha_j \quad (\text{bzw. } f_j \ge \hat f_j - \alpha_j \text{ für die Maximierungen } j = 7, 8, 11)
  \qquad \forall j < k, \label{obj:lex}
\end{equation}
wobei $\hat f_j$ der in Stufe $j$ erreichte Wert ist. Gurobi setzt das als Mehrzieloptimierung mit fallenden
Prioritäten um (\code{setObjectiveN}). Die Toleranzen $\alpha_j$ (\code{abstol}) sind:
\begin{center}
\small
\begin{tabular}{l l l}
\hline
Stufe & Gesamtspielraum $\varepsilon_j$ (\code{stage\_slack}) & $\alpha_j = \varepsilon_j/2$ \\
\hline
$f_4$ \code{gap\_spread} & $2$ & $1$ \\
$f_5$ \code{staffing\_spread} & $|U^{D}| \cdot n / 7$ & $|U^{D}| \cdot n / 14$ \\
$f_9$ \code{shift\_balance} & $2\,|U^{D}| \cdot n / 7$ & $|U^{D}| \cdot n / 7$ \\
alle anderen & -- & $10^{-6}$ \\
\hline
\end{tabular}
\end{center}

\textbf{Erklärung.}
Eine spätere Stufe darf eine frühere nur um deren Toleranz verschlechtern. Bei den drei Verteilungsstufen $f_4$,
$f_5$ und $f_9$ ist das bewusst etwas Spielraum, damit die späteren Stufen noch Wahlmöglichkeiten haben; die
andere Hälfte des Spielraums erlaubt einen frühen Abbruch der Stufe (siehe unten). Insgesamt liegt der Endwert
einer solchen Stufe höchstens $\varepsilon_j$ über dem bestmöglichen Wert. Alle anderen Stufen halten ihren Wert
praktisch exakt.

\textbf{Abbruchkriterien.}
\begin{center}
\small
\begin{tabular}{l l}
\hline
Kriterium & Wert \\
\hline
Gesamtzeit $T$ (\code{SOLVER\_MAX\_TIME\_SECONDS}) & in der Demo fest 30\,s \\
relative Optimalitätslücke (\code{MIPGap}) & 3\,\%; exakt (0) für $f_1$, $f_2$, $f_3$, $f_6$ und $f_8$ \\
absolute Optimalitätslücke (\code{MIPGapAbs}) & $\varepsilon_j/2$ für $f_4$, $f_5$, $f_9$ \\
Zeitlimit einzelner Stufen (\code{stage\_time\_share}) & $0{,}15\,T$ für $f_5$, $0{,}10\,T$ für $f_6$, $f_8$ und $f_9$ \\
Watchdog & beendet die Rechnung spätestens 5\,s nach $T$ \\
\hline
\end{tabular}
\end{center}

\textbf{Erklärung.}
$f_1$ und $f_3$ werden exakt gelöst: 3\,\% eines großen Werts (ein unerreichbares Konto oder viele Lücken) würden
weitere Konten außerhalb des Bands oder vermeidbare Lücken zulassen; ebenso die kleine Zahl $f_2$ der geteilten
Wochenenden am Monatsanfang. Auch $f_6$ wird exakt gelöst, weil schon ein Samstag mehr oder weniger zählt, und
$f_8$, weil der Wert in die Hunderte geht und 3\,\% davon eine Person in vielen Fenstern ohne Nachtdienst lassen
könnten. Die Zeitlimits schützen die folgenden Stufen: Bei $f_5$, $f_6$, $f_8$ und $f_9$ kann die Schranke den
erreichbaren Wert verfehlen (Bruchteile von Personen, durch Abwesenheiten oder die Wochenendregeln unerreichbare
Werte), sodass Gurobi die Stufe sonst bis zur Gesamtzeit weiterrechnen würde. Die übrigen Stufen laufen, bis sie bewiesen sind oder $T$ erreicht
ist. Ist $T$ erreicht, gilt der beste bis dahin gefundene Plan; spätere Stufen kommen dann nicht mehr an die
Reihe.

%==========================================================================================================
\section{Nach dem Lösen}

\textbf{Ergebnis.} Aus der Lösung werden die Dienste $\{(i,s,d,u) : x_{isdu} = 1\}$ mit der Qualifikation
$\lambda_{iud}$ übernommen, nur für Tage des Monats: Ein mitgeplanter Sonntag gehört nicht zum Plan. Die
gemeldeten Lücken werden nicht aus $g$ gelesen, sondern aus den Diensten berechnet, $\max(0,\, r_{sdul} - \sum_{i
\in I_{udl}} x_{isdu})$, ebenso die angezeigten Werte von $f_3$, $f_4$, $f_5$, $f_6$, $f_8$ und $f_9$. So stimmen
sie genau mit der unabhängigen Prüfung überein, auch wenn Hilfsvariablen wegen der Toleranzen \eqref{obj:lex} über
ihrem echten Wert liegen.

\textbf{Unabhängige Prüfung.} Die Prüfung (\code{api/app/domain/acceptance.py}) übernimmt vom Solver nichts außer
den Diensten. Sie rechnet mit den echten Uhrzeiten und den Diensten des Vormonats und prüft unter anderem:
Mindestbesetzung und gemeldete Lücken, Zugehörigkeit und Qualifikation, Einschränkungen samt keiner Nacht vor einem
Abwesenheitstag und keinem Spätdienst vor Urlaubsbeginn oder Fortbildung \eqref{eq:NN}, \eqref{eq:NS}, einen Dienst
pro Tag,
den Monatssaldo innerhalb $\pm\tau$, tägliche Höchstarbeitszeit, Pausen und den Durchschnitt von 8\,h je Werktag,
11\,h Ruhezeit, höchstens $N^{max}$ Nächte in Folge und 48\,h Erholung nach der letzten Nacht. Dazu kommen die
Regeln des Hauses: höchstens $K^{S}$ Arbeitstage am Stück (Fortbildung zählt mit), am Wochenende dieselbe Schicht an
beiden Tagen und nie zwei gearbeitete Wochenenden in Folge, jeweils auch über den Monatswechsel, sofern der
Vormonat geladen ist. Ein Wochenende über den Monatswechsel, das \eqref{obj:weekend} nicht fortführen konnte, wird
also abgelehnt. Abweichend vom Gesetzestext gilt in der Demo-Version:
\begin{itemize}
  \item \textbf{Spät $\to$ Früh} am Folgetag ist von der 11-h-Ruhezeit ausgenommen. Das Modell hält diese Wechsel
        über \eqref{obj:ltoe} klein, die Prüfseite zählt sie als Kennzahl. Jede andere Ruhezeit unter 11\,h wird
        abgelehnt.
  \item \textbf{Erholung nach Nächten}: Auch eine Fortbildung an einem der $L^{N}$ Tage nach der letzten Nacht
        ist ein Verstoß, passend zu \eqref{con:recoverytraining}. Ein Nachtdienst-Wunsch in den $L^{N}$ Tagen vor
        einer Fortbildung gilt deshalb als nicht erfüllbar, ebenso nach \eqref{eq:NS} ein Spätdienst-Wunsch am Tag
        davor.
  \item \textbf{Zeitumstellung}: Modell und Prüfung ignorieren sie. Eine Schicht zählt für die tägliche
        Höchstarbeitszeit und die Pausen mit ihren geplanten Zeiten, auch die Nacht der Umstellung im Oktober,
        die tatsächlich eine Stunde länger dauert.
  \item \textbf{Ersatzruhetag}: Ob jeder gearbeitete Sonn- oder Feiertag durch einen freien Werktag ausgeglichen
        wird (ArbZG §11 Abs.~3), prüft die Demo-Version nicht; das Modell bildet diese Regel ebenfalls nicht ab.
\end{itemize}

\end{document}
